% Moved from appendix_proofs.tex to the supplement (page budget).
% ------------------------------------------------------------
\section{Separable Dissimilarities Admit Exact Near-Linear-Time Evaluation}
\label{app:separable}

Balanced designs are needed only when $d$ admits no algebraic shortcut;
the following proposition gives two such shortcuts.

\begin{proposition}[Separable dissimilarities]
\label[proposition]{prop:separable}
Let $D$ be a dissimilarity table.
\begin{enumerate}
\item[(i)] Suppose $d_{ij}=f_i+f_j-2\langle\phi_i,\phi_j\rangle$ for scalars $f_i$
and vectors $\phi_i$ in a $q$-dimensional inner-product space. Then
\[
\SND(D)=\frac1N\Big[(n-1)\sum_i f_i+\sum_i\norm{\phi_i}^2-\Big\|\sum_i\phi_i\Big\|^2\Big],
\]
which is computable exactly in $O(nq)$ operations. Examples include squared
Euclidean distance ($f_i=\norm{x_i}^2$, $\phi_i=x_i$), giving the variance
identity ${\SND=2\sum_i\norm{x_i-\bar x}^2/(n-1)}$. Further
examples are squared $W_2$ between Gaussians $\mathcal N(m_i,\Sigma_i)$ with
pairwise commuting (for example diagonal) covariances, with
$f_i=\norm{m_i}^2+\operatorname{tr}\Sigma_i$ and $\phi_i=(m_i,\Sigma_i^{1/2})$; squared
Hellinger distance $\tfrac12\lVert\sqrt{\pi_i}-\sqrt{\pi_j}\rVert^2$ between
distributions $\pi_i$ on $q$ outcomes, with $\phi_i=\sqrt{\pi_i}/\sqrt2$ and $f_i=\frac12$;
cosine dissimilarity $1-\cos(x_i,x_j)$, with $\phi_i=x_i/(\sqrt2\,\norm{x_i})$ and
$f_i=\frac12$; LPIPS-type squared feature distances; and $0$--$1$
disagreement between categorical outputs ($\phi_i=e_{y_i}/\sqrt2$, $f_i=\frac12$).
\item[(ii)] Suppose $d_{ij}=\sum_{c=1}^q\abs{x_{ic}-x_{jc}}$, an $\ell_1$
dissimilarity. This includes total variation between categorical action
distributions, averaged over observations. Then
\[
\sum_{i<j}d_{ij}=\sum_{c=1}^q\sum_{t=1}^n(2t-n-1)\,x_{(t)c},
\]
where $x_{(1)c}\le\dots\le x_{(n)c}$ are the sorted values of coordinate $c$. This is
computable exactly in $O(q\,n\log n)$ operations.
\end{enumerate}
\end{proposition}

\begin{proof}
\emph{(i)} $\sum_{i<j}(f_i+f_j)=(n-1)\sum_if_i$, and
$\sum_{i<j}\langle\phi_i,\phi_j\rangle=\frac12\big(\norm{\sum_i\phi_i}^2-\sum_i\norm{\phi_i}^2\big)$.
For squared Euclidean distance, $(n-1)\sum_i\norm{x_i}^2+\sum_i\norm{x_i}^2-\norm{\sum_i x_i}^2=n\sum_i\norm{x_i-\bar x}^2$,
and $N=n(n-1)/2$. \emph{(ii)} For sorted values,
$\sum_{s<t}(x_{(t)}-x_{(s)})=\sum_t x_{(t)}\big[(t-1)-(n-t)\big]$, and the
coordinates add.
\end{proof}

\begin{remark}[Non-separable dissimilarities and metric scope]
Many dissimilarities are of neither form (i) nor form (ii). Examples are unsquared
Euclidean distance in dimension $q\ge2$, Hellinger distance, $W_2$ between
Gaussians (hence SND as defined by \citealt{bettini2025snd}), Jensen--Shannon
distance, edit and optimal-transport distances between trajectories, and
rollout- or judge-based comparisons. Univariate $W_1$, by contrast, is the
$\ell_1$ distance between cumulative distribution functions, so for
distributions on a common finite grid it is covered by (ii). Distances with explicit Euclidean
embeddings, $d_{ij}=\norm{x_i-x_j}_2$, admit a further, randomized shortcut. Since
$\E_g\abs{\langle v,g\rangle}=\sqrt{2/\pi}\,\norm v$ for $g\sim\mathcal N(0,I_q)$, random
one-dimensional projections combined with (ii) give unbiased estimates of $\mu$
that touch every individual. This is the projection-based approach of
\citet{barhum2007approximating}, and its accuracy is set by the number of
projections. Balanced designs need no embedding, only black-box access to $d$,
and are therefore the tool of choice when no embedding is available.
Theorem~\ref{thm:balanced} covers every dissimilarity table, whether or not it is
a metric. Theorem~\ref{thm:metric} applies only to pseudometric tables. These
include exactly computed Euclidean, Hellinger, Jensen--Shannon, edit and
$W_p$ optimal-transport distances (including $W_2$); they exclude comparisons that can violate the
triangle inequality: Kullback--Leibler divergences, cross-play return gaps, entropic (Sinkhorn)
transport costs, dynamic-time-warping alignments, judge scores, and Monte Carlo
estimates of distances that use different samples for different pairs (averages over one shared sample, such as the common observation pool of SND, remain pseudometrics).
\end{remark}

